% !PRINTED 114
\weberpage{114}

Hiernach führen wir die drei Funktionen ein:
\weq{9}{
\begin{aligned}
f(\omega) &= \frac{e^{-\frac{\pi i}{24}}\eta\left(\dfrac{\omega+1}{2}\right)}{\eta(\omega)},\\
f_1(\omega) &= \frac{\eta\left(\dfrac{\omega}{2}\right)}{\eta(\omega)},\\
f_2(\omega) &= \sqrt{2}\,\frac{\eta(2\omega)}{\eta(\omega)}.
\end{aligned}
}

Dann ist nach (6), (7), (8)
\weq{10}{
\begin{aligned}
\vartheta_{00} &= \eta(\omega)f(\omega)^2,\\
\vartheta_{01} &= \eta(\omega)f_1(\omega)^2,\\
\vartheta_{10} &= \eta(\omega)f_2(\omega)^2,
\end{aligned}
}
und aus (8) und (9) folgt, daß $f(\omega)$, $f_1(\omega)$,
$f_2(\omega)$ für ein rein imaginäres $\omega$ reell und positiv
sind. (10) stimmt überein mit § 24, (10), und die dort
eingeführten Funktionen $f$, $f_1$, $f_2$ sind dieselben wie
diese. Aus (10) folgen aber leicht die Relationen [(2) und
§ 21 (14)]:
\weq{11}{
\begin{aligned}
f(\omega)^8 &= f_1(\omega)^8 + f_2(\omega)^8,\\
\sqrt{2} &= f(\omega)f_1(\omega)f_2(\omega).
\end{aligned}
}

Mit Hilfe dieser Formeln kann man durch Quadratwurzeln jede der
drei Funktionen $f(\omega)^8$, $f_1(\omega)^8$, $f_2(\omega)^8$
durch jede andere ausdrücken. Dazu führt die aus (11) fließende
Formel:
\weq{12}{
\begin{aligned}
{[f_1(\omega)^8 - f_2(\omega)^8]^2} &= f(\omega)^{16} - \frac{64}{f(\omega)^8},\\
{[f(\omega)^8+f_2(\omega)^8]^2} &= f_1(\omega)^{16} + \frac{64}{f_1(\omega)^8},\\
{[f(\omega)^8+f_1(\omega)^8]^2} &= f_2(\omega)^{16} + \frac{64}{f_2(\omega)^8}.
\end{aligned}
}

Für die linearen Fundamentaltransformationen der Funktionen $f$
erhält man zunächst aus (4) und (9):
\weq{13}{
\begin{aligned}
f(\omega+1) &= e^{-\frac{\pi i}{24}}f_1(\omega),\\
f_1(\omega+1) &= e^{-\frac{\pi i}{24}}f(\omega),\\
f_2(\omega+1) &= e^{\frac{\pi i}{12}}f_2(\omega).
\end{aligned}
}
